%% LyX 2.3.4.2 created this file.  For more info, see http://www.lyx.org/.
%% Do not edit unless you really know what you are doing.
\documentclass[11pt,spanish]{article}
\renewcommand{\rmdefault}{cmr}
\usepackage[T1]{fontenc}
\usepackage[utf8]{inputenc}
\usepackage{geometry}
\geometry{verbose,tmargin=2cm,bmargin=2cm,lmargin=2cm,rmargin=2cm}
\setlength{\parskip}{\bigskipamount}
\setlength{\parindent}{0pt}
\usepackage{amsthm}
\usepackage{amssymb}

\makeatletter

%%%%%%%%%%%%%%%%%%%%%%%%%%%%%% LyX specific LaTeX commands.
%% Because html converters don't know tabularnewline
\providecommand{\tabularnewline}{\\}

%%%%%%%%%%%%%%%%%%%%%%%%%%%%%% Textclass specific LaTeX commands.
\theoremstyle{remark}
\newtheorem{rem}{\protect\remarkname}[section]
\theoremstyle{definition}
\newtheorem{xca}{\protect\exercisename}[section]
\theoremstyle{definition}
\newtheorem{example}{\protect\examplename}[section]

%%%%%%%%%%%%%%%%%%%%%%%%%%%%%% User specified LaTeX commands.
\usepackage{titling}
\setlength{\droptitle}{-2cm}
\usepackage{multicol}
\usepackage{enumitem}
\usepackage{proof}
\usepackage{mathrsfs}
\usepackage{amsfonts}
\usepackage{forest}
\usepackage{ebproof}
\usepackage{amssymb}
\definecolor{Guinda}{rgb}{0.49, 0.05, 0.12}
\definecolor{Rojo}{rgb}{0.8, 0, 0}
\definecolor{Verde}{rgb}{0,0.4,0}
\definecolor{Azul}{rgb}{0,0.4,0.4}
\definecolor{Naranja}{rgb}{0.8,0.4,0}
\usepackage{tikz}
\usetikzlibrary{automata, positioning, arrows}

\makeatother

\usepackage{babel}
\addto\shorthandsspanish{\spanishdeactivate{~<>.}}

\usepackage{listings}
\lstset{mathescape=true}
\providecommand{\examplename}{Ejemplo}
\providecommand{\exercisename}{Ejercicio}
\providecommand{\remarkname}{Observación}

\begin{document}
\title{Teoría de la Computación, 2021-1\\
Nota de clase 4: Expresiones Regulares}
\author{Manuel Soto Romero}
\date{\today \\ Colegio de Ciencia y Tecnología UACM}

\maketitle
\renewcommand*{\thefootnote}{\arabic{footnote}}
\setcounter{footnote}{0}
\setcounter{section}{4}
\renewcommand*\lstlistingname{C\'odigo}

Continuaremos nuestro estudio de los Lenguajes Regulares por medio
de un método basado en métodos algebraicos conocidos como \emph{expresiones
regulares}. Estas expresiones regulares (como veremos más adelante)
pueden definir exactamente los mismos lenguajes que describen las
gramáticas regulares y los autómatas finitos que estudiamos con anterioridad:
lenguajes regulares. 

Muchos sistemas como lenguajes de programación, procesadores de textos,
sistemas gestores de bases de datos y muchos otros, incluyen este
mecanismo como un método de procesamiento de cadenas. Como ejemplos
de programas tenemos al comando \texttt{grep} que permite buscar
texto en los sistemas operativos basados en \textsc{Unix} o el software
para construir analizadores léxicos \textsc{JFlex} de gran importancia
en el diseño e implementación de compiladores. En esta nota nos dedicaremos
a estudiar este tipo de mecanismos.

\subsection{Definición}

Sea $\Sigma$ un alfabeto. La \emph{expresión regular} sobre $\Sigma$
y los conjuntos que ellas denotan son definidos recursivamente como
sigue:
\begin{itemize}
\item $\varnothing$ es una expresión regular y denota el conjunto vacío.
\item $\varepsilon$ es una expresión regular y denota el conjunto $\left\{ \varepsilon\right\} $.
\item Para cada $a$ en $\Sigma$, $a$ es una expresión regular y denota
el conjunto $\left\{ a\right\} $.
\item Si $r$ y $s$ son expresiones regulares denotando el lenguaje $R$
y $S$ respectivamente entonces $\left(r+s\right),rs,r^{*}$ son expresiones
regulares que denotan los conjuntos $R\cup S,RS$ y $R^{*}$, respectivamente.
\end{itemize}
De esta forma, usamos a las expresiones regulares como una simplificación
para denotar lenguajes. Recordando que los lenguajes son conjuntos
podemos contrastar la definición anterior en la siguiente tabla:
\begin{center}
\begin{tabular}{|c|c|}
\hline 
Expresión Regular & Lenguaje\tabularnewline
\hline 
$\varnothing$ & $\varnothing$\tabularnewline
\hline 
$a$ & $\left\{ a\right\} $\tabularnewline
\hline 
$\varepsilon$ & $\left\{ \varepsilon\right\} $\tabularnewline
\hline 
$\left(r+s\right)$ & $R\cup S$\tabularnewline
\hline 
$\left(rs\right)$ & $RS$\tabularnewline
\hline 
$\left(r^{*}\right)$ & $R^{*}$\tabularnewline
\hline 
\end{tabular}
\par\end{center}

Los paréntesis sólo son agrupadores. Además, se tiene la siguiente
prioridad y asociatividad de operadores:
\begin{itemize}
\item La cerradura $*$ asocia a la izquierda y tiene la mayor prioridad.
\item Posteriormente se realiza la concatenación que asocia a la derecha.
\item Finalmente se tiene el operador $+$ de alternativa que también asocia
por la derecha.
\end{itemize}
Por ejemplo:

\[
0+1b^{*}+a1^{*}
\]

se resuelve en el siguiente orden:

\[
(0+((1(b^{*}))+(a\left(1^{*}\right))))
\]

Veamos algunos ejemplos de expresiones regulares junto con el lenguaje
que reconocen.
\begin{itemize}
\item ¿Qué lenguaje reconoce la siguiente expresión regular?

\[
\left(0+1\right)^{*}
\]

Analicemos cada subexpresión:
\begin{itemize}
\item $0$ es la expresión regular que denota al lenguaje $\left\{ 0\right\} $
\item $1$ es la expresión regular que denota al lenguaje $\left\{ 1\right\} $
\item Entonces $\left(0+1\right)$ es el lenguaje $\left\{ 0\right\} \cup\left\{ 1\right\} =\left\{ 0,1\right\} $
\item De esta forma $\left(0+1\right)^{*}=\left\{ 0,1\right\} ^{*}=\left\{ \varepsilon,0,1,00,01,10,11,000,001,...\right\} $
\end{itemize}
Es decir, reconoce el lenguaje de todas las posibles cadenas que se
pueden formar como 0 y 1.

\[
L_{1}=\left\{ w:w\in\Sigma^{*}\right\} 
\]

\item ¿Qué lenguaje reconoce la siguiente expresión regular?

\[
\left(a+b\right)c^{*}
\]

Analicemos cada subexpresión:
\begin{itemize}
\item $a=\left\{ a\right\} $
\item $b=\left\{ b\right\} $
\item $c=\left\{ c\right\} $
\item $\left(a+b\right)=\left\{ a,b\right\} $
\item $c^{*}=\left\{ \varepsilon,c,cc,ccc,cccc,...\right\} $
\item $\left(a+b\right)c^{*}=\left\{ a,b,ac,bc,acc,bcc,...\right\} $
\end{itemize}
El lenguaje admitido es el de todas las cadenas que comienzan con
$a$ o $b$ seguidas de 0 o más símbolos $c$.
\end{itemize}
\bigskip{}

\begin{rem}
La representación de lenguajes regulares por medio de expresiones
regulares no es única. Es posible que haya varias expresiones regulares
para el mismo lenguaje. Por ejemplo, $b\left(a+b\right)^{*}$ y $b\left(b+a\right)^{*}$
representan el mismo lenguaje. Otro ejemplo sería con las expresiones
$\left(a+b\right)^{*}$ y $\left(a^{*}b^{*}\right)^{*}$.
\end{rem}
\bigskip{}

\begin{xca}
Indica el lenguaje que reconocen las siguientes expresiones regulares:
\end{xca}
\begin{enumerate}
\item $\left(a^{*}c\right)+\left(ab^{*}\right)$
\item $\left(00\right)^{*}+\left(1\left(11\right)^{*}\right)$
\end{enumerate}
Veamos ahora, algunos ejemplos sobre diseño de expresiones regulares. 

\bigskip{}

\begin{example}
Encontrar expresiones regulares que representen los siguientes lenguajes,
definidos sobre el alfabeto $\Sigma=\left\{ a,b\right\} $:
\end{example}
\begin{enumerate}
\item Lenguaje de todas las cadenas que comienzan con el símbolo $b$ y
terminan con el símbolo $a$.
\begin{itemize}
\item Debemos garantizar que las cadenas inicien con $b$, lo cual puede
ser representado por: $b$
\item Después del primer símbolo, puede haber cualquier número de $a$s
y $b$s, de forma tal que podemos generar dichas cadenas con la expresión:
$\left(a+b\right)^{*}$.
\item Finalmente, la cadena debe terminar con $a$, con lo cual, la expresión
regular más idonea es: $a$.
\item Concatenando las tres expresiones anteriores tenemos:

\emph{Solución: $b\left(a+b\right)^{*}a$}
\end{itemize}
\item Lenguaje de todas las cadenas que tienen un número par de símbolos.
\begin{itemize}
\item Dado que queremos pares de símbolos necesitamos construir cada uno
de éstos, es decir, tanto el primer como el segundo símbolo pueden
ser una $a$ o una $b$ y debemos concatenarlos. Esto se puede representar
con $\left(a+b\right)\left(a+b\right).$
\item Con la expresión del punto anterior, estamos construyendo cadenas
de longitud 2, pero podemos tener otras longitudes, por ejemplo 0,
2, 4, 6, etc. Dado que ya concatenamos dos símbolos, esto se puede
lograr aplicando la estrella de Kleene. Es decir:

\emph{Solución: $\left(\left(a+b\right)\left(a+b\right)\right)^{*}$}
\end{itemize}
\item Lenguaje de todas las cadenas que tienen un número par de $a$s.
\begin{itemize}
\item Para reconocer cadenas que contengan un número par de $as$, partamos
de un caso fácil, la cadena que tiene 2 de estos símbolos. En este
caso, podemos tener alguna de las siguientes formas:
\begin{itemize}
\item Dos símbolos al principio.
\item Dos símbolos al final.
\item Dos símbolos en medio.
\item Dos símbolos en cualquier posición.
\end{itemize}
Podríamos considerar todos estos casos y separarlos con una operación
de suma, sin embargo, podemos auxiliarnos de la estrella de Kleene
para indicar que puede haber cero o más presencias del símbolo $b$
entre las dos $a$ que queremos. Es decir: $b^{*}ab^{*}ab^{*}$.
\item Ahora que tenemos fijado un par de $a$s, podemos generalizarlo como
hicimos en el ejemplo anterior con una estrella de Kleene. Es decir:
$\left(b^{*}ab^{*}ab^{*}\right)^{*}$.
\item Hasta el punto anterior, pareciera que ya tenemos la expresión regular,
sin embargo, recordemos que el cero también es un número par, de forma
que una cadena sin $a$s también pertenece al lenguaje. Podemos agregar
este paso con una suma, quedando:

\emph{Solución: $\left(b^{*}ab^{*}ab^{*}\right)$$^{*}+b^{*}$}
\end{itemize}
\end{enumerate}
\bigskip{}

\begin{example}
Encontrar expresiones regulares que representen los siguientes lenguajes,
definidos sobre el alfabeto $\Sigma=\left\{ 0,1\right\} $
\end{example}
\begin{enumerate}
\item Lenguaje de todas las cadenas que tienen exactamente dos ceros.
\begin{itemize}
\item Este es un caso particular de los ejemplos que vimos anteriormente.
Dado que queremos exactamente dos ceros, simplemente colocamos los
unos correspondientes con su respectiva estrella de Kleene entre los
dos símbolos $0$:

\emph{Solución: $1^{*}01^{*}01^{*}$}
\end{itemize}
\item Lenguaje de todas las cadenas cuyo penúltimo símbolo, de izquierda
a derecha es un 0.
\begin{itemize}
\item Queremos que el penúltimo símbolo sea un 0, de forma que antes de
encontrar dicho símbolo, podemos tener cualquier número de ceros y
unos, es decir: $\left(0+1\right)^{*}$
\item Posteriormente, aparecerá el símbolo 0, que se reconoce con: $0$
\item Y finalmente, después de dichos patrones, podemos tener ya sea otro
cero u otro uno, es decir: $\left(0+1\right)$
\item Con lo cual, concatenando las 3 expresiones anteriores, tenemos:

\emph{Solución: $\left(0+1\right)^{*}0\left(0+1\right)$}
\end{itemize}
\end{enumerate}
\bigskip{}

\begin{xca}
Sea $\Sigma=\left\{ 0,1\right\} $. Encontrar una expresión regular
que represente el lenguaje de todas las cadenas que no tienen dos
ceros consecutivos.
\end{xca}
%
\begin{xca}
Sea $\Sigma=\left\{ a,b,c\right\} $. Encontrar una expresión regular
que represente el lenguaje de todas las cadenas que no contiene la
cadena $bc$.
\end{xca}
%

\begin{thebibliography}{1}
\bibitem{key-1}Hopcroft, J. E; Montwani, R., and Ullman, J. \emph{Introduction
to Automata Theory Language, and Computation}. Addison Wesley, third
edition 2007.

\bibitem{key-1}Rodrigo De Castro, \emph{Notas de Clase de Teoría
de la Computación: Lenguajes, autómatas, gramáticas,} Universidad
Nacional de Colombia, Facultas de Ciencias, 2004.
\end{thebibliography}

\end{document}
